% Coverage certificate for heart-s4-c1 (genome c85804fa-08bd-4207-ac69-a31a8567473c); content sha256 5bb264e6720c1a25b94714c430f44b3b1f28f41592eef373c0a2d422a381b5a0.
\providecommand{\explaneatcertdir}[1]{}
\begin{drafted}[heart-s4-c1 coverage certificate]
\section{Coverage certificate: heart-s4-c1}\label{app:heart-s4-c1:sec}
Provenance, coverage and checks of the explanation, rendered from the canonical JSON record \texttt{heart-s4-c1.json} (schema 1.4.0, SHA-256 \texttt{5bb264e6720c1a25}\ldots). The explanation itself (summary, parts, evidence, figures) is the read-down document before this page.

\begin{table}[htbp]
\centering\small
\begin{tabularx}{\linewidth}{@{}lX@{}}
\toprule
Dataset & Heart Disease (binary, prepared) \\
Experiment & 10seed-heart-s4, seed 4 \\
Genome & \texttt{c85804fa-08bd-4207-ac69-a31a8567473c} (NEAT key 9679, generation 39) \\
Explanation & \texttt{10565a81-f55b-47f4-b51e-21641ec08556} (36 operations) \\
Train AUC / accuracy & 0.9423 / 0.8884 (n = 242) \\
Test AUC / accuracy & 0.8858 / 0.8197 (n = 61) \\
Output activation & sigmoid \\
Phenotype & 31 nodes (26 inputs, 4 hidden, 1 output), 29 enabled edges \\
Explained model state & 36 nodes, 29 enabled edges \\
Refactoring operations & split\_\allowbreak{}node 5 \\
Function preservation & max $|\Delta\,\mathrm{output}|$ = $3.33\times 10^{-16}$ on the training split \\
Annotations & 5 (4 leaf, 1 composition) \\
Structural coverage & 1 (35 of 35 candidate nodes; 6 Null) \\
Compositional coverage & 1 \\
Uncovered nodes / edges & 0 / 0 \\
Missing compositions & 0 \\
Full coverage & yes \\
Drafted by & AI-assisted draft \\
Verified by & unverified \\
Record hash (SHA-256) & \texttt{5bb264e6720c1a25}\ldots \\
\bottomrule
\end{tabularx}
\caption{Explanation certificate summary: heart-s4-c1.}
\label{app:heart-s4-c1:summary}
\end{table}
\subsection{Coverage}\label{app:heart-s4-c1:coverage}
Coverage is computed by the current code (the function behind the \texttt{get\_coverage} tool) under these conventions. \emph{Candidate nodes} are all nodes of the explained model state except the output nodes. A candidate node is \emph{covered} when some annotation of the hierarchy (a leaf, or a composition together with its descendants) contains it and all its enabled outgoing edges. \emph{Exact Null} parts, which have no effect on any output for any input, count as covered: nodes with no path to an output over enabled, non-zero-weight edges (disconnected inputs and dead ends), and edges that have zero weight or touch such a node. Compositional coverage is the fraction of required compositions present: every annotation except a single global root containing every output must have a parent; a parentless, childless annotation made only of Null nodes is exempt.

\begin{itemize}
\item \textbf{Null, disconnected input (6):} \texttt{-12} (restecg:0), \texttt{-14} (restecg:2), \texttt{-16} (exang:0), \texttt{-24} (thal:4), \texttt{-26} (thal:7), \texttt{-5} (cp:2).
\item \textbf{Uncovered nodes:} none.
\item \textbf{Uncovered edges:} none.
\item \textbf{Missing compositions:} none.
\end{itemize}
\noindent\textbf{Agreement with the \texttt{get\_coverage} tool:} yes.

\subsection{Closed-form checks}\label{app:heart-s4-c1:closedform}
{\small\noindent\begin{tabularx}{\linewidth}{@{}lXrr@{}}
\toprule Tag & Annotation & Closed form vs forward pass & Exact function vs forward pass \\ \midrule
A1 & heart\_\allowbreak{}disease\_\allowbreak{}logit & 0.00419 & $2.78\times 10^{-16}$ \\
A2 & fixed\_\allowbreak{}defect\_\allowbreak{}normal\_\allowbreak{}sugar\_\allowbreak{}gate & $4.34\times 10^{-4}$ & 0 \\
A3 & no\_\allowbreak{}asymptomatic\_\allowbreak{}pain\_\allowbreak{}indicator & 0.00173 & 0 \\
A4 & downsloping\_\allowbreak{}st\_\allowbreak{}indicator & $1.95\times 10^{-4}$ & 0 \\
A5 & thalach\_\allowbreak{}floor\_\allowbreak{}relay & $8.60\times 10^{-4}$ & 0 \\
\bottomrule
\end{tabularx}\par}
\subsection{Refactoring operations}\label{app:heart-s4-c1:operations}
The explained model state is the phenotype after these operations, replayed in order. All are function-preserving; the largest change in the model output over the 242 training rows is $3.33\times 10^{-16}$ (tolerance $1.00\times 10^{-9}$).
\begin{itemize}
\item \#26: split \texttt{-10} into \texttt{-10\_\allowbreak{}a}, \texttt{-10\_\allowbreak{}b}.
\item \#27: split \texttt{-11} into \texttt{-11\_\allowbreak{}a}, \texttt{-11\_\allowbreak{}b}.
\item \#28: split \texttt{-15} into \texttt{-15\_\allowbreak{}a}, \texttt{-15\_\allowbreak{}b}.
\item \#29: split \texttt{-21} into \texttt{-21\_\allowbreak{}a}, \texttt{-21\_\allowbreak{}b}.
\item \#30: split \texttt{-7} into \texttt{-7\_\allowbreak{}a}, \texttt{-7\_\allowbreak{}b}.
\end{itemize}
\subsection{Audit against the definitions}\label{app:heart-s4-c1:audit}
Automatic findings on the stored explanation (codes as in the JSON record's \texttt{issues}); they flag what an examiner should check, not verdicts.
\begin{itemize}
\item \textsc{warning} \texttt{unverified}: No human verification is recorded for this explanation.
\end{itemize}
\end{drafted}
